% ====================================================================
\section{P-completeness of program evaluation}\label{sec:pcomplete}
% ====================================================================

\begin{definition}\label{def:eval}
$\EVAL$ is the problem of deciding, for a given MSLP $\mathcal P$, whether $\val(\mathcal P)=1$.
$\OMCVP$ (topologically ordered monotone circuit value) is the restriction of $\EVAL$ to programs
in normal form.
\end{definition}

The following theorem is classical.

\begin{theorem}[Ladner; Goldschlager]\label{thm:mcvp}
$\OMCVP$ is $\Ptime$-complete under logspace many-one reductions.
\end{theorem}

\begin{proof}[Sources]
Ladner~\cite{Lad75} proved $\Ptime$-completeness of the circuit value problem. In his encoding a
circuit is a straight-line sequence of assignments, each referring only to earlier entries.
Goldschlager~\cite{Gol77} showed that the restriction to monotone circuits (constants, $\AND$, $\OR$)
remains $\Ptime$-complete under logspace reductions. His double-rail transformation replaces each
gate by a constant number of monotone gates that refer only to the images of earlier gates, so it
preserves the topologically ordered presentation. Appending the gate $\AND(o,o)$ preserves the
answer and is a logspace operation. For a uniform treatment of these variants see
\cite[Appendix~A.1]{GHR95}. The source notes of the original manuscript use the same ordered
convention~\cite{Src}.
\end{proof}

\begin{corollary}\label{cor:evalcomplete}
$\EVAL$ is $\Ptime$-complete under logspace many-one reductions.
\end{corollary}

\begin{proof}
Membership: evaluate the instructions in order, which takes $O(L)$ operations on a table of $r$
bits. Hardness: $\OMCVP$ is the special case of programs in normal form.
\end{proof}

\begin{remark}[Two different space bounds]\label{rem:twospaces}
Two space bounds play different roles in the rest of the paper, and they must not be confused.
\begin{itemize}
\item The \emph{reduction} of \cref{thm:simulation} runs in logarithmic space. It reads the program
  text and writes the tree. It never evaluates the program and never stores register values.
  Everything it outputs is determined by the instruction types and operand indices: vertices,
  edges, initial bits, control words, and the query time.
\item \emph{Evaluating} the program is $\Ptime$-complete by \cref{cor:evalcomplete}. It is not
  known to be possible in logarithmic space, which would imply $\Lsp=\Ptime$. In the constructed
  tree, the evaluation is carried out by the majority dynamics itself, during the physical time
  interval $[0,\tau]$.
\end{itemize}
Obliviousness (\cref{sec:model}) is what separates the two. The schedule is a function of the
program text only, so the transducer can write it without knowing any computed bit.
\end{remark}
